\section{\texorpdfstring{The semi-norms \(\mathbf{N}_p\)}{The semi-norms N sub p}}

Let \(p\) be a real number \(\geqslant 1\) and let \(\mathcal{F}^p(\mathrm{X},M)\) be the set of finite numerical functions \(f\) defined on \(\mathrm{X}\) such that \(M(|f|^p)\) is finite. It is obvious that if \(g\) is a function belonging to \(\mathcal{F}^p(\mathrm{X},M)\) and if \(|f| \leqslant |g|\) , then \(f\) also belongs to \(\mathcal{F}^p(\mathrm{X},M)\) ; this remark and Minkowski's inequality show that the sum of two functions in \(\mathcal{F}^p(\mathrm{X},M)\) also belongs to this set; taking into account the fact that \(M\) is positively homogeneous, we thus see that \(\mathcal{F}^p(\mathrm{X},M)\) is a linear subspace of the space \(\mathbf{R}^{\mathrm{X}}\) of all finite numerical functions defined on \(\mathrm{X}\) .

For every number \(p > 0\) and every finite numerical function \(f\) defined on \(X\) , set

\[
\mathrm{N} _ {p} (f) = \left(M (| f | ^ {p})\right) ^ {1 / p};
\]

then \(\mathrm{N}_p(\lambda f) = |\lambda |\mathrm{N}_p(f)\) for every scalar \(\lambda\) ; moreover, if \(p \geqslant 1\) then, by (7),

\[
\mathrm{N} _ {p} (f + g) \leqslant \mathrm{N} _ {p} (f) + \mathrm{N} _ {p} (g),\tag{8}
\]

which proves that \(N_{p}\) is a semi-norm on the vector space \(\mathcal{F}^{p}(X,M)\) (TVS, II, §1).

PROPOSITION 4. --- Let p and q be two real numbers \textgreater{} 0 and set \(1/r = 1/p + 1/q\) . For any finite numerical functions f, g defined on X,

\[
\mathrm{N} _ {r} (f g) \leqslant \mathrm{N} _ {p} (f) \mathrm{N} _ {q} (g)\tag{9}
\]

provided that \(\mathrm{N}_p(f)\) and \(\mathrm{N}_q(g)\) are finite.

For, the relation (9) may be written

\[
M (| f | ^ {r} | g | ^ {r}) \leqslant \bigl (M (| f | ^ {p}) \bigr) ^ {r / p} \bigl (M (| g | ^ {q}) \bigr) ^ {r / q},
\]

which is none other than Hölder's inequality (5) applied to the numbers \(\alpha = r / p\) and \(\beta = r / q\) and to the functions \(|f|^p\) and \(|g|^q\) .

COROLLARY. --- Suppose that \(M(1) = 1\) ; then, for every finite numerical function f defined on X, the mapping \(p \mapsto \mathrm{N}_{p}(f)\) is increasing in {]}0, +∞{[}.

Applying the inequality (9) to the case that g = 1, we see that \(\mathrm{N}_{r}(f) \leqslant \mathrm{N}_{p}(f)\) for all q \textgreater{} 0; since the number r defined by \(1/r = 1/p + 1/q\) runs over the set of numbers such that 0 < r < p when q runs over the set of numbers \textgreater{} 0, the corollary is proved.

PROPOSITION 5. --- For every finite numerical function f defined on X, the set I of values of 1/p (p \textgreater{} 0) such that \(\mathrm{N}_{p}(f)\) is finite is either empty or is an interval; if I is not reduced to a point, then the mapping \(\alpha \mapsto \log \mathrm{N}_{1/\alpha}(f)\) is either convex on I or is equal to \(-\infty\) on the interior of I.

Let r and s be two distinct numbers \textgreater{} 0 such that 1/r and 1/s belong to I; it all comes down to proving that if

\[
{\frac {1}{p}} = {\frac {t}{r}} + {\frac {1 - t}{s}},
\]

with \(0 < t < 1\) , then

\[
\log \mathrm{N} _ {p} (f) \leqslant t \cdot \log \mathrm{N} _ {r} (f) + (1 - t) \log \mathrm{N} _ {s} (f),\tag{10}
\]

or, what comes to the same,

\[
\mathrm{N} _ {p} (f) \leqslant \left(\mathrm{N} _ {r} (f)\right) ^ {t} \left(\mathrm{N} _ {s} (f)\right) ^ {1 - t},\tag{11}
\]

a relation that may, by the definition of \(N_{p}\) , be written

\[
M (| f | ^ {p}) \leqslant \left(M (| f | ^ {r})\right) ^ {t p / r} \left(M (| f | ^ {s})\right) ^ {(1 - t) p / s}.\tag{12}
\]

Setting \(\alpha = tp / r\) , we have \(1 - \alpha = (1 - t)p / s\) by the relation that defines \(p\) as a function of \(t, r, s\) ; whence \(p = \alpha r + (1 - \alpha)s\) . Hölder's inequality now yields

\[
M \big (| f | ^ {r \alpha} | f | ^ {s (1 - \alpha)} \big) \leqslant \big (M (| f | ^ {r}) \big) ^ {\alpha} \big (M (| f | ^ {s}) \big) ^ {1 - \alpha},
\]

which is none other than the inequality (12).

